\documentclass[a4paper,11pt]{jsarticle}

% --- パッケージ ---
\usepackage[top=25mm,bottom=25mm,left=25mm,right=25mm]{geometry}
\usepackage{amsmath}
\usepackage{booktabs}
\usepackage{tabularx}
\usepackage{url}
\usepackage{enumitem}
\usepackage{graphicx}
\usepackage{listings}
\usepackage{color}

\lstset{
  basicstyle=\ttfamily\small,
  frame=single,
  breaklines=true,
  columns=fullflexible
}

\title{LLM駆動型避難シミュレーションの構築と\\交通ネットワーク再現における課題}
\author{}
\date{}

\begin{document}

\maketitle

% =============================================================================
\section{はじめに}
% =============================================================================

本稿では，大規模言語モデル（LLM）を意思決定エンジンとするマルチエージェント避難シミュレーションの構築過程を報告するとともに，現時点で未解決の技術的課題である交通シミュレーションの導入に関する考察を行う．

本シミュレーションは，東日本大震災における福島県沿岸部（いわき市・浪江町・南相馬市）の避難行動を対象としている．
国土交通省のProject PLATEAUが提供する3D都市モデルをUnityゲームエンジン上に展開し，最大500体のLLM駆動型エージェントが各自のペルソナ（年齢，身体能力，家族構成，災害経験等）に基づいて自律的に避難判断を下す．
エージェントの移動にはUnityのNavMeshを用いた歩行者移動を採用しており，避難所選択，家族捜索，周辺住民との会話，他者への追従といった社会的行動を含む6種のアクションを実装している．

現在のシステムは歩行者移動のみを対象としており，車両避難時の渋滞やそれに伴う避難者の心理的変化を再現できていない．
東日本大震災の避難実態調査では自動車による避難割合が地域によっては80\%を超えており\cite{evacuationsurvey}，交通状況の再現は避難シミュレーションの信頼性を左右する重要な要素である．
本稿では，まずシミュレーション環境の構築過程で遭遇した技術的問題とその解決を記録し，次に交通シミュレーション導入の障壁となっている道路データの構造的制約を分析する．

% =============================================================================
\section{シミュレーション環境の構築}
% =============================================================================

本シミュレーションの環境構築には，Unity，PLATEAU SDK，Pythonベースの LLMサーバ等，複数の技術スタックの統合が必要である．
以下ではリポジトリの取得からシミュレーション実行に至るまでの手順と，構築過程で遭遇した主要な技術的問題を記述する．

\subsection{前提条件}

実行環境として以下が必要である（表\ref{tab:requirements}）．

\begin{table}[htbp]
\centering
\caption{動作に必要なソフトウェア環境}
\label{tab:requirements}
\begin{tabularx}{\textwidth}{lXl}
\toprule
ソフトウェア & 要件 & 備考 \\
\midrule
Unity & 2023.2.19f1 & URP 16.0.6 を使用 \\
Python & 3.10 以上 & パッケージマネージャに uv を使用 \\
OpenAI API キー & GPT-4o-mini & LLMエージェント機能に必須 \\
\bottomrule
\end{tabularx}
\end{table}

レンダーパイプラインとしてUniversal Render Pipeline（URP）16.0.6を採用しているため，ビルトインレンダーパイプラインでは動作しない．
この制約は後述するマテリアル互換性の問題に直結する．

\subsection{セットアップ手順}

リポジトリのクローンからシミュレーション実行までの主要な手順は以下の通りである．

\subsubsection{PLATEAU SDKの配置}

PLATEAU SDK（\texttt{com.synesthesias.plateau-unity-sdk}）はファイルサイズの関係からGit管理対象外としており，手動でのインストールが必要である．
PLATEAU SDK for Unity\cite{plateausdk}のリリースページから\texttt{.tgz}ファイルをダウンロードし，\texttt{Packages/com.synesthesias.plateau-unity-sdk/}に展開する．

SDKの参照方式は開発過程で複数回変更されている．
当初は\texttt{manifest.json}でローカルの絶対パスを指定していたが環境依存性の問題が生じ，GitHub URL参照（\texttt{git\#v2.3.2}）への変更を経て，最終的にGit管理外のローカルパッケージとして手動管理する方式に落ち着いた．

\subsubsection{LLMサーバの設定}

LLMエージェント機能はPythonのWebSocketサーバを介してOpenAI APIと通信する構成をとっている．
\texttt{llm\_server/}ディレクトリ内で環境変数（APIキー，モデル名，サーバアドレス・ポート）を設定し，\texttt{uv sync}で依存関係を解決する．
Unityシーン上の\texttt{LLMServerManager}コンポーネントにより，再生開始時にサーバプロセスが自動起動される．

主な設定項目を表\ref{tab:llmconfig}に示す．

\begin{table}[htbp]
\centering
\caption{LLMサーバの設定項目}
\label{tab:llmconfig}
\begin{tabularx}{\textwidth}{lll}
\toprule
環境変数 & デフォルト値 & 説明 \\
\midrule
\texttt{OPENAI\_API\_KEY} & （必須） & OpenAI APIキー \\
\texttt{OPENAI\_MODEL} & \texttt{gpt-4o-mini} & 使用するLLMモデル \\
\texttt{LLM\_SERVER\_HOST} & \texttt{127.0.0.1} & WebSocketサーバアドレス \\
\texttt{LLM\_SERVER\_PORT} & \texttt{8765} & WebSocketサーバポート \\
\bottomrule
\end{tabularx}
\end{table}

\subsubsection{シーンファイルの入手}

シミュレーション対象地域のシーンファイル（いわき市，浪江町，南相馬市）はそれぞれ100\,MBを超えるため，Git管理外としている．
PLATEAUの3D都市モデルデータをUnityにインポートし，NavMeshのBake等の前処理を施した状態のシーンファイルを別途共有する運用としている．

\subsection{主要パッケージ構成}

本プロジェクトで使用する主要なUnityパッケージを表\ref{tab:packages}に示す．

\begin{table}[htbp]
\centering
\caption{主要Unityパッケージ一覧}
\label{tab:packages}
\begin{tabularx}{\textwidth}{lll}
\toprule
パッケージ & バージョン & 用途 \\
\midrule
\texttt{com.synesthesias.plateau-unity-sdk} & ローカル管理 & 3D都市モデル統合 \\
\texttt{com.unity.ai.navigation} & 2.0.6 & NavMeshによる経路探索 \\
\texttt{com.unity.ml-agents} & 2.0.2 & 強化学習エージェント \\
\texttt{com.unity.render-pipelines.universal} & 16.0.6 & URPレンダリング \\
\texttt{com.unity.inputsystem} & 1.7.0 & 新入力システム \\
\bottomrule
\end{tabularx}
\end{table}

\subsection{構築過程で遭遇した技術的問題}

シミュレーション環境の構築にあたり，複数の技術的問題に遭遇した．
以下ではそのうち主要なものを記述する．

\subsubsection{マテリアル互換性の問題}

PLATEAUの3D都市モデルをUnityにインポートした際，建物モデルがピンク色（シェーダーエラー）で表示される現象が発生した．
原因は，PLATEAU SDKがビルトインレンダーパイプライン向けのマテリアルを生成するのに対し，本プロジェクトではURPを採用していたことにある．
URP用のRender Pipeline Assetの作成，\texttt{GraphicsSettings}および\texttt{QualitySettings}の更新，ならびに全マテリアルのURP対応シェーダーへの変換により解決した．

\subsubsection{PLATEAU RoadNetworkのコンパイルエラー}

PLATEAU SDK Toolkitが\texttt{PLATEAU.Editor.RoadNetwork.RoadNetworkEditMode}型を参照していたが，使用していたSDKバージョンにはこの型が存在せず，コンパイルエラーが発生した．
暫定的な対処として，\texttt{None}，\texttt{Edit}，\texttt{View}の3値を持つenum型のスタブを作成し，コンパイルを通す措置をとった．
この問題は，PLATEAU SDKとToolkitのバージョン間の互換性管理が十分でないことに起因する．

\subsubsection{NavMeshにおける道路の地形埋没}

PLATEAUの道路メッシュ（\texttt{tran}オブジェクト群）がDEM（地形モデル）より低い位置に配置されているケースがあり，NavMesh上での経路探索で避難者エージェントが地面に埋まった状態で移動する現象が生じた．
これに対し，道路メッシュの各頂点から地形に向けてレイキャストを行い，地形表面の高さにクリアランスを加えた位置に道路メッシュを自動調整するエディタスクリプト（\texttt{RoadSnapper}）を作成した．
上下双方向のレイキャストにより，地形より上・下いずれのケースにも対応している．

\subsubsection{エピソードリセットの不備}

複数エピソードを連続実行する実験において，前エピソードの動的状態（会話履歴，家族の合流状態，行動履歴等）がリセットされず，LLMの応答異常やエージェントの不正な振る舞いが発生した．
全動的状態を初期化する\texttt{ResetForNewEpisode()}メソッドを追加し，エピソード切替時にこれを呼び出すよう修正した．

\subsubsection{ログ出力の不完全性}

シミュレーション実行時のアクションログおよび避難率CSV等が正しく出力されない，または不完全な状態で終了する問題が発生した．
ログ出力ロジックの改修，記録タイミングの修正，およびバッチ実験用設定機能の追加により対応した．

% =============================================================================
\section{交通シミュレーション導入の課題}
% =============================================================================

本シミュレーションにおける最も重要な未解決課題の一つが，車両避難を含む交通シミュレーションの導入である．
以下では，この課題を3D都市モデルの道路データの構造的制約，現在のシステムアーキテクチャとの整合性，および要件定義の観点から分析する．

\subsection{3D都市モデルにおける道路データの制約}

\subsubsection{LODレベルと交通シミュレーションの関係}

PLATEAUの3D都市モデルでは，道路データのLOD（Level of Detail）に応じて利用可能な情報が大きく異なる（表\ref{tab:lod}）．

\begin{table}[htbp]
\centering
\caption{LODレベルごとの道路データ特性}
\label{tab:lod}
\begin{tabularx}{\textwidth}{lXXXX}
\toprule
LOD & 形状表現 & 車線情報 & 接続情報 & 交通シミュレーション \\
\midrule
LOD1 & ポリゴン（面） & なし & なし & 不可 \\
LOD2 & ポリゴン＋道路種別 & 一部あり & なし & 不可 \\
LOD3 & 詳細な道路モデル & 車線定義あり & 接続関係あり & 可能 \\
\bottomrule
\end{tabularx}
\end{table}

PLATEAUのSandbox Toolkitが提供する交通シミュレーション機能はLOD3道路データを前提としている\cite{plateautraffic}．
本シミュレーションの対象地域である福島県内では，いわき市が2020年度のLOD1データを，南相馬市が2022年度のLOD1を主としたデータを提供しているにとどまる．
浪江町についてはPLATEAUの3D都市モデル自体が提供されていない可能性がある\cite{plateauportal}．
したがって，PLATEAUのToolkitによる交通シミュレーションを対象地域に適用することは現時点では困難である．

なお，PLATEAU公式ツール群にはLOD1から道路ネットワーク（ノード・リンクデータ）を生成可能なPLATEAU-RoadNetwork-Generator\cite{roadnetgen}が存在するが，LOD1入力時の出力品質（特に交差点接続の成功率やトポロジの正確さ）に関する定量的な評価は公開されておらず，実用性は未検証である．

\subsubsection{道路データの孤立点構造}

現在のシステムでは，PLATEAUの道路オブジェクト（CityGMLの\texttt{tran\_}プレフィクスを持つオブジェクト群）を空間インデックスとして事前にBake（編集時の事前計算）し，実行時の検索に供している．
各道路オブジェクトから抽出される情報を表\ref{tab:roadentry}に示す．

\begin{table}[htbp]
\centering
\caption{道路データ（RoadEntry）の保持属性}
\label{tab:roadentry}
\begin{tabularx}{\textwidth}{llX}
\toprule
属性 & 型 & 内容 \\
\midrule
\texttt{name} & 文字列 & 道路オブジェクト名 \\
\texttt{position} & Vector3 & バウンディングボックス中心座標 \\
\texttt{functionCode} & 整数 & 道路機能コード（1=高速，2=国道，3=県道，4=市町村道） \\
\texttt{sectionType} & 整数 & 区間種別（1=土工，2=高架橋，3=橋梁，4=交差部，5=アンダーパス，6=トンネル） \\
\texttt{gmlId} & 文字列 & PLATEAU GML ID \\
\bottomrule
\end{tabularx}
\end{table}

各道路オブジェクトの情報はバウンディングボックスの中心座標1点と分類属性のみであり，道路の線形（頂点列），隣接道路への参照，方向，車線数，幅員といったネットワーク構築に不可欠な情報は一切保持されていない．
この結果，実行時の道路検索は「ある地点から半径$r$メートル以内に存在する道路の一覧」を返す半径ベースの点検索に限定され，「ある道路の先にどの道路が接続しているか」というグラフ探索は構造的に不可能である．

\subsubsection{道路ポリゴンの分断問題}

PLATEAUのLOD1道路データは，Unityシーン上では個別のGameObject（個別メッシュ）として存在する．
これらのメッシュは物理的に隣接していても論理的な接続情報を持たない．具体的には，(1)~隣接するポリゴンが同一道路の一部なのか別の道路なのか判別できない，(2)~交差点ポリゴン（区間種別コード4）においてどの道路ポリゴンが合流しているかの情報がない，(3)~メッシュの端が微妙にズレたり重なったりする「つぎはぎ問題」が存在する，という3点が挙げられる．

前節で述べた道路メッシュの高さ調整（\texttt{RoadSnapper}による地形スナップ処理）は垂直方向の位置補正のみであり，水平方向における道路間の接続性には寄与しない．
すなわち，LOD1データのみでは道路ポリゴンの集合をネットワーク（ノードとエッジのグラフ）へ変換する処理が本質的に困難であり，これが交通シミュレーション導入の根本的な障壁となっている．

\subsection{システムアーキテクチャとの整合性}

\subsubsection{NavMeshとの共存問題}

現在のシステムでは全避難者エージェントおよび緊急車両エージェント（消防車巡回）がUnityのNavMeshAgentを用いて移動している．
NavMeshは「歩行可能領域」として一枚の連続メッシュ上で経路計算を行うものであり，道路・歩道・広場の区別を内在しない．

車両避難者を導入する場合，NavMesh上で速度パラメータのみを変更する方法では車両が歩道や建物間の隙間を通過する非現実的な挙動が生じる．
一方，交通ネットワークグラフ上で別個に移動させる方法では，二つの全く異なる経路探索システムの並存による設計の複雑化が避けられない．
特に，渋滞に巻き込まれた車両避難者が車を放棄して徒歩に切り替えるシナリオ---東日本大震災で実際に多数発生した事象\cite{evacuationsurvey}---の再現には，交通ネットワーク上の位置からNavMesh上の位置へのハンドオフ処理が必要となり，二つのシステム間の座標・状態の受け渡しが技術的課題となる．

\subsubsection{LLMへの情報伝達の不在}

現在のLLMエージェントに提供される環境コンテキストには，最寄りの主要道路名と道路種別の2項目のみが含まれており，交通量・渋滞状況・道路の通行可否といった動的な交通情報は一切伝達されていない．
交通シミュレーションを導入したとしても，その結果をLLMの意思決定に反映させるためには，交通状況を適切な粒度で自然言語化し，エージェントの判断コンテキストに組み込む情報伝達パイプラインの新規構築が必要である．

\subsubsection{計算性能への影響}

現在のシステムでは500体のLLMエージェントが各ステップでWebSocket経由のLLM API呼び出しを行っており，建物・道路データはグリッドベース空間インデックス（セルサイズ50\,m）により$O(1)$に近い検索性能を実現している．
交通シミュレーションの追加に伴い，グラフ上の経路探索（ダイクストラ法等），道路区間ごとの車両数のリアルタイム集計，渋滞の動的更新と伝播計算といった処理が加わる．
NavMeshの経路計算がUnityエンジン内部で高度に最適化されているのに対し，交通ネットワーク上のグラフ探索は自前実装となるため，500エージェント規模での実行性能への影響は慎重に検証する必要がある．

\subsection{要件定義上の課題}

交通シミュレーションの導入にあたっては，道路データやアーキテクチャの技術的制約に加え，シミュレーション設計として以下の要件が現時点で未定義である．

\begin{description}[leftmargin=2em, style=nextline]
  \item[避難手段の選択基準]
  避難者のうち車両避難と徒歩避難の比率をどのように決定するか．現在のペルソナデータには車両保有情報が含まれておらず，家族構成（高齢者・幼児の同行）が手段選択に与える影響も定式化されていない．

  \item[車両の空間的配置]
  避難者が車両にアクセスするまでの過程（駐車場の位置，車両への到達時間，乗車時間）および車両台数の上限・相乗り避難の仕組みが未設計である．

  \item[災害による道路の損壊・閉塞]
  地震による道路損壊，がれきによる通行止め，液状化による道路陥没，津波到達に伴う通行不可区間の動的生成，橋梁損壊による孤立地域の発生といった，災害時特有の道路状態変化の表現方法が未定義である．
  橋梁に関しては区間種別コード3として既にデータ上の識別は可能であるが，損壊の発生条件やそのネットワークへの反映方法は定められていない．

  \item[避難手段の動的切替]
  渋滞に巻き込まれた車両避難者が車を放棄して徒歩に切り替えるシナリオの実装方法が未定義である．
  放棄された車両が道路を塞ぐことによる二次的渋滞の連鎖的再現も課題である．

  \item[検証指標の特定]
  交通シミュレーションの妥当性を評価するための定量的指標が定められていない．
  東日本大震災時の車両プローブデータ（Honda・Toyotaが研究用に提供した走行実績データ）や国土交通省の道路交通センサスとの比較検証が候補として考えられるが，対象地域・対象災害に完全に合致するデータの取得は容易でない．
\end{description}

% =============================================================================
\section{おわりに}
% =============================================================================

本稿では，LLM駆動型マルチエージェント避難シミュレーションの構築過程を記述するとともに，交通シミュレーション導入に関する技術的課題を整理した．

環境構築においては，PLATEAU SDKの参照方式の安定化，URPとのマテリアル互換性の確保，道路メッシュの地形埋没への対処，エピソードリセットの整備など，3D都市モデルとゲームエンジンの統合に伴う固有の問題群に対応した．
これらの知見は，PLATEAUの3D都市モデルをUnity上でリアルタイムシミュレーションに活用する際の実践的な参考となりうる．

交通シミュレーションについては，対象地域の道路データがLOD1に限定されていることが根本的な制約である．
LOD1データでは車線情報・接続情報が欠落しており，道路ポリゴンの集合からネットワーク構造（ノードとエッジのグラフ）を自動的に構築することが本質的に困難である．
加えて，NavMeshに基づく既存の移動システムとの共存，LLMへの交通情報伝達パイプラインの不在，500エージェント規模での計算性能，避難手段の選択基準や道路損壊表現等の要件定義の不足といった課題が重層的に存在する．

これらの課題に対しては，外部道路ネットワークデータ（OpenStreetMap，国土数値情報等）の活用，LOD1ポリゴンからの自動ネットワーク生成，および交通の抽象化モデルの3つのアプローチが考えられるが，いずれも単独では十分でなく，段階的な検証と組み合わせが必要である．
今後の取り組みとして，まず最小限の交通情報をLLMエージェントに提供し，意思決定への影響を定量的に計測した上で，交通シミュレーションの精緻化の方向性を判断することが合理的と考える．

% =============================================================================
% 参考文献
% =============================================================================
\begin{thebibliography}{99}
\bibitem{evacuationsurvey}
内閣府, ``東日本大震災時の地震・津波避難に関する住民アンケート調査,'' 2012.

\bibitem{plateautraffic}
Project PLATEAU, ``PLATEAU SDK Toolkits for Unity --- Sandbox Toolkit,'' \\
\url{https://github.com/Project-PLATEAU/PLATEAU-SDK-Toolkits-for-Unity}, 参照 2026-03-05.

\bibitem{plateausdk}
Project PLATEAU, ``PLATEAU SDK for Unity,'' \\
\url{https://github.com/Project-PLATEAU/PLATEAU-SDK-for-Unity}, 参照 2026-03-05.

\bibitem{roadnetgen}
Project PLATEAU, ``PLATEAU-RoadNetwork-Generator,'' \\
\url{https://github.com/Project-PLATEAU/PLATEAU-RoadNetwork-Generator}, 参照 2026-03-05.

\bibitem{plateauportal}
国土交通省, ``PLATEAUオープンデータポータル,'' \\
\url{https://www.mlit.go.jp/plateau/open-data/}, 参照 2026-03-05.
\end{thebibliography}

\end{document}
